\documentclass[12pt]{article}
\usepackage[utf8]{inputenc}
\usepackage[T1]{fontenc}
\usepackage[margin=1in]{geometry}
\usepackage{setspace}
\usepackage{xcolor}
\usepackage[colorlinks=true, urlcolor=blue!60!black, linkcolor=blue!60!black]{hyperref}
\setlength{\parindent}{2em}
\setlength{\parskip}{1.25em}
\renewcommand{\baselinestretch}{1.0}
\begin{document}
\begin{center}
{\LARGE \textit{Seila Law LLC v. Consumer Financial Protection Bureau}, 591 U.S. 197 (2020)\par}
\vspace{0.8em}
{\large Case Brief --- Four-Page Edition\par}
\end{center}
\vspace{1em}

\section{Memory Jogger}

Congress violates the separation of powers when it vests significant executive authority in an independent agency headed by a single Director removable only for cause, though the removal restriction is severable from the rest of the Dodd-Frank Act (\href{https://www.law.cornell.edu/supremecourt/text/19-7}{Cornell LII}).

\section{Detailed Case Facts}

After the 2008 financial crisis, the Dodd-Frank Act created the Consumer Financial Protection Bureau, giving it administration of 18 existing consumer-protection statutes, a new ban on unfair, deceptive, or abusive practices, and power to make rules, issue civil investigative demands, adjudicate, and sue for civil penalties. Departing from the multimember-commission model, Congress placed the CFPB under a single Director serving a five-year term, removable only for ``inefficiency, neglect of duty, or malfeasance in office,'' and funded it from the Federal Reserve rather than annual appropriations. In 2017 the Bureau served a civil investigative demand on Seila Law LLC, a debt-relief law firm. The firm objected that the CFPB's structure was unconstitutional; the Bureau ordered compliance.

\section{Procedural History}

The Central District of California enforced the demand, and the Ninth Circuit affirmed, holding that \textit{Humphrey's Executor} and \textit{Morrison v. Olson} controlled. In the Supreme Court, the Solicitor General and Director Kraninger agreed that the removal restriction was unconstitutional but argued that it was severable and that Acting Director Mulvaney had ratified the demand. The Court appointed Paul Clement as amicus to defend the judgment.

\section{Judicial Votes}

Vacated and remanded (decided June 29, 2020).
\begin{itemize}
\item \textbf{Removal (5--4):} Chief Justice Roberts, joined by Justices Thomas, Alito, Gorsuch, and Kavanaugh; Justice Kagan, joined by Justices Ginsburg, Breyer, and Sotomayor, dissenting.
\item \textbf{Severability (7--2):} Roberts, joined by Alito and Kavanaugh (plurality); Kagan's four concurring in the judgment; Justice Thomas, joined by Justice Gorsuch, dissenting.
\end{itemize}

\section{Holding}

Article II vests the entire executive power in the President, which generally includes power to remove executive officers at will. Precedent recognizes only two exceptions: multimember, partisan-balanced expert agencies that do not wield substantial executive power (\textit{Humphrey's Executor}) and inferior officers with limited duties (\textit{Morrison}). For-cause protection for a single Director wielding significant executive power is unconstitutional, but under Dodd-Frank's severability clause it is severed, leaving the CFPB in operation. The Court did not overrule \textit{Humphrey's Executor}, treated the funding and term as aggravating features only, and left ratification for remand.

\section{Analysis of Opinions}

\textbf{Majority (Roberts).} Seila Law had standing, the claim was ripe, and the case remained adverse because the Government still sought enforcement. On the merits, the Vesting and Take Care Clauses, the Decision of 1789, \textit{Myers}, and \textit{Free Enterprise Fund} make removal the rule. The CFPB fits neither exception: its Director is a single principal officer making rules under 19 statutes, adjudicating, and seeking penalties. The structure lacks historical footing, its supposed predecessors being brief or contested, and it clashes with a design that divides power everywhere except the Presidency, where a single President is accountable to the whole electorate. A five-year term and funding outside appropriations aggravated the problem.

\textbf{Plurality on severability (Roberts).} Under \textit{Alaska Airlines} and the express severability clause, the rest of the statute is fully operative, and abolishing the agency would cause great disruption; Congress may redesign the CFPB as a commission if it wishes.

\textbf{Thomas (concurring in part, dissenting in part).} The Court should overrule \textit{Humphrey's Executor}, whose quasi-legislative and quasi-judicial powers are a fiction creating a fourth branch. Courts have only the negative power to disregard unconstitutional laws, not to excise them; he would simply deny enforcement of the demand.

\textbf{Kagan (dissenting in part).} The text says nothing about removal; the Necessary and Proper Clause lets Congress structure offices so long as the President can perform his duties. History, from the Comptroller of the Treasury to the Federal Reserve, shows insulated financial regulators. A single Director is \emph{easier} for a President to oversee than a gaggle of staggered commissioners.

\section{Examples: Future Applications}

\begin{enumerate}
\item \textbf{Same-Side (Invalid): FHFA Director.} A single FHFA Director removable only for cause wields significant executive power---the result in \textit{Collins v. Yellen} (2021).
\item \textbf{Same-Side (Invalid): Independent EPA Administrator.} A single Administrator with a six-year protected term who issues nationwide rules and seeks penalties.
\item \textbf{Opposite-Side (Valid): Bipartisan Maritime Commission.} Five staggered, partisan-balanced commissioners who only adjudicate fee disputes fit \textit{Humphrey's Executor}.
\item \textbf{Opposite-Side (Valid): Special Inspector.} An inferior officer auditing one grant program and referring fraud fits \textit{Morrison}.
\item \textbf{Fence-Sitter: The Modern FTC or SEC.} Multimember boards that now wield substantial enforcement power. \textit{Seila Law} preserved \textit{Humphrey's Executor} but narrowed it; in \textit{Trump v. Slaughter} (2026) the Court invalidated the FTC's protection.
\end{enumerate}

\section{Critique}

\textbf{Originalist views.} Unitary-executive scholars such as Bamzai and Prakash defend the premise but fault Roberts for stopping halfway. Revisionists---Shugerman, Manners and Menand, Katz and Rosenblum---argue that the Decision of 1789 settled nothing, early Congresses regulated tenure, and the for-cause language had specific historical meanings.

\textbf{Progressive views.} Functionalists such as Peter Shane see judicial aggrandizement that ignores Congress's power to insulate financial regulators from capture, and call the single-head versus multimember line arbitrary. The case pits accountability and formalism against expertise and legislative flexibility.

\section{Key Quotations}

\begin{quote}
``Under our Constitution, the `executive Power'---all of it---is `vested in a President,' who must `take Care that the Laws be faithfully executed.' Art. II, \S 1, cl. 1; id., \S 3.''
\end{quote}
\noindent --- Chief Justice Roberts, for the Court.

\begin{quote}
``The CFPB's single-Director structure contravenes this carefully calibrated system by vesting significant governmental power in the hands of a single individual accountable to no one.''
\end{quote}
\noindent --- Chief Justice Roberts, for the Court.

\begin{quote}
``We think it clear that Congress would prefer that we use a scalpel rather than a bulldozer in curing the constitutional defect we identify today.''
\end{quote}
\noindent --- Chief Justice Roberts, plurality opinion (Part IV).

\end{document}
